% !PRINTED 233
\weberpage{233}
%<<CONT
restrict ourselves in this proof to the consideration of the icosahedral group, since for the other groups entirely similar inferences are to be made, which we may leave to the reader. The circumstance that in the tetrahedral and octahedral groups relative invariants occur alongside the absolute ones, while the icosahedral group, as a simple group, has only absolute invariants, is of no essential influence here.

We prove, then, that all invariants of the icosahedral group can be represented as integral rational functions of the three forms $f$, $H$, $T$.

To this we are led by the following chain of inferences:

1. \gesperrt{No two of the icosahedral forms $f$, $H$, $T$ have a common divisor.}

This follows immediately from the definition of these functions, according to which the equations $f = 0$, $H = 0$, $T = 0$ furnish the five-fold, three-fold and two-fold poles, which are all different from one another.

2. \gesperrt{A double root of the icosahedral equation}
\weq{1}{H^3 - z f^5 = 0}
\gesperrt{is necessarily a root of $f$, $H$ or $T$, and can therefore occur only for one of the values $z = 0, \infty, 1728$.}

For if (1) has a double root, then together with the function the first derivative, or, if one employs the homogeneous form, by \gesperrt{Euler}'s theory [Bd. I, \S.\,17, (5)] the two derivatives with respect to $x_1$ and $x_2$, must vanish at the same time, hence
\[
3H^2 H'(x_1) - 5z f^4 f'(x_1) = 0,
\]
\[
3H^2 H'(x_2) - 5z f^4 f'(x_2) = 0;
\]
consequently, since by 1. $H$ and $f$ do not vanish at the same time, only three possibilities remain: either $H = 0$, $z = 0$, or $f = 0$, $z = \infty$, or finally
\[
H'(x_1)\,f'(x_2) - H'(x_2)\,f'(x_1) = 0,
\]
i.e. $T = 0$, $z = 1728$ [\S.\,60, (4)]. This also yields the evident theorem:
%CONT>>
